

\chapter{Einleitung und Problemhintergrund}\label{cha:introduction}

Im Zuge von Industrie 4.0 werden Maschinen und Produktionsanlagen zunehmend vernetzt. Sensoren erfassen laufend ihren Betriebszustand, während Cloud-Infrastrukturen die entstehenden Daten verarbeiten. Predictive Maintenance nutzt diese Daten, um Veränderungen früh zu erkennen und Wartungen besser zu planen. Dadurch können ungeplante Stillstände und Instandhaltungskosten reduziert werden \parencite[1--2, 7]{zontaPredictiveMaintenanceIndustry2020}.

Ein Teil davon ist die Prognose der \ac{RUL}. Sie beschreibt die geschätzte Zeit beziehungsweise Anzahl an Betriebszyklen, bis eine Maschine oder Komponente eine festgelegte Ausfallgrenze erreicht. \ac{RUL}-Prognosen helfen somit bei der Entscheidung, wann eine Wartung notwendig ist \parencite[2]{huRemainingUsefulLife2019}.

Für die datenbasierte \ac{RUL}-Prognose werden unter anderem \ac{LSTM}, \acp{GRU}, \acp{CNN} und \acp{TCN} eingesetzt. Diese Modelle erlernen anhand historischer Sensordaten zeitliche Zusammenhänge und Degradationsmuster, aus denen sie die verbleibende Nutzungsdauer ableiten \parencites[15--16]{wuRemainingUsefulLife2024}[3, 6]{wangRemainingUsefulLife2020}.

Ein verbreiteter Benchmark ist der mit \ac{CMAPSS} erzeugte Datensatz der \ac{NASA}. Er enthält simulierte Sensordaten von Turbofan-Triebwerken, deren Trainingsverläufe bis zum Ausfall reichen. Die ursprünglichen Messreihen enthalten keine fehlenden Sensorwerte. In realen Anwendungen können dagegen Sensor- und Kommunikationsfehler einzelne Messwerte oder ganze Sensorkanäle ausfallen lassen. Auch systematische Messabweichungen und Sensordrift können auftreten. Solche Störungen können die \ac{RUL}-Prognose verschlechtern und werden daher bereits durch gezielte Veränderungen der Benchmarkdaten untersucht \parencite[7--10]{tangDeepLearningFramework2026}.

Neben den etablierten Modellen werden inzwischen \acp{TSFM} untersucht. Modelle wie MOMENT werden auf umfangreichen Zeitreihensammlungen aus unterschiedlichen Anwendungsbereichen vortrainiert. Die dabei erlernten Merkmale können anschließend für weitere Aufgaben verwendet werden \parencite[1--4]{goswamiMOMENTFamilyOpen2024}. Dintén und Zorrilla nutzen die von MOMENT erzeugten Merkmale bereits für die \ac{RUL}-Prognose auf \ac{CMAPSS}, untersuchen jedoch keine beschädigten Testdaten \parencite[246, 252--255]{dintenUsingTimeSeries2025}. El-Ghoussani et al. kombinieren wiederum Chronos-2 mit einem kleinen Regressionskopf und nennen die Robustheit bei fehlenden Sensoren als offene Forschungsfrage \parencite[2--4]{el-ghoussaniTimeSeriesFoundationModel2026}.

Die geplante Arbeit vergleicht deshalb MOMENT-1-small mit einem \ac{LSTM} und einem \ac{TCN} unter gezielt veränderten Testbedingungen. Bei MOMENT bleibt der vortrainierte Encoder unverändert, während ein zusätzlicher Regressionskopf für die \ac{RUL}-Schätzung trainiert wird. Fehlende Messwerte, Übertragungslücken, Sensorausfälle und Drift sollen zeigen, wie empfindlich die Modelle auf eine nachlassende Datenqualität reagieren. Neben der Prognosegüte werden ihr Ressourcenbedarf und die Kosten des Cloud-Betriebs untersucht. Damit soll beurteilt werden, welchen Nutzen der gewählte \ac{TSFM}-Ansatz gegenüber den beiden für diese Aufgabe trainierten Modellen bietet.
